% --- START OF FILE 05_ui_liste.tex ---

\chapter{UI e Liste: Da ListView a Jetpack Compose}

La costruzione di interfacce grafiche in Android ha subito un'enorme evoluzione.
È vitale per l'esame conoscere sia l'approccio classico (XML e View) che quello moderno dichiarativo (Jetpack Compose).

\section{Approccio Classico: View e ViewGroup}
Il sistema classico si basa sull'albero degli oggetti UI, un'implementazione del \textbf{Design Pattern Composite}.
\begin{itemize}
    \item \textbf{\texttt{View}} (Component): La classe base per tutti gli elementi grafici interattivi (\texttt{TextView}, \texttt{Button}, ecc.). Implementa la funzione \texttt{onDraw()}.
    \item \textbf{\texttt{ViewGroup}} (Composite): Sottoclasse di \texttt{View} che funge da contenitore (Layout) per altre View o ViewGroup (es. \texttt{LinearLayout}, \texttt{FrameLayout}).
\end{itemize}

\subsection{Layout in XML}
Tra i \texttt{ViewGroup} più usati:
\begin{itemize}
    \item \textbf{\texttt{LinearLayout}:} Dispone gli elementi in una singola riga o colonna.
    Usa la proprietà \texttt{android:layout\_weight} per definire in proporzione \textit{quanto spazio extra} deve occupare un figlio rispetto agli altri.
    \item \textbf{\texttt{ConstraintLayout}:} Layout piatto (flat), basato su vincoli (constraint) tra gli elementi.
    Molto performante per UI complesse rispetto a LinearLayout innestati.
    \item \textbf{\texttt{FrameLayout}:} Progettato per contenere un singolo elemento principale.
    Se ha più figli, li \textbf{sovrappone} uno sopra l'altro (funziona da "cornice").
\end{itemize}

\section{Le Liste nel Sistema Classico: ListView e Adapters}
Per mostrare dati dinamici (es.
rubrica contatti) non si possono creare migliaia di TextView manualmente.
Si usano i componenti \texttt{AdapterView} (\texttt{ListView}, \texttt{GridView}) accompagnati dagli \textbf{Adapter}.

L'\textbf{Adapter} è una classe non-visiva che funge da "ponte" tra l'origine dei dati (Array, Database) e la \texttt{ListView}.
I principali sono \texttt{ArrayAdapter}, \texttt{CursorAdapter} e il generico \texttt{BaseAdapter}.

\subsection{Il concetto di View Recycling (Domanda da Esame!)}
Creare nuovi oggetti in Java/Kotlin (\textit{Object creation}) è un'operazione "costosa" in termini di memoria e CPU, e sollecita troppo il \textit{Garbage Collector}.
Per questo motivo, la \texttt{ListView} instanzia in RAM \textbf{solo il numero di righe strettamente necessarie} a riempire lo schermo, più un paio di scorta.

Quando l'utente fa scroll verso il basso:
\begin{enumerate}
    \item La riga che esce dallo schermo in alto finisce in una coda di riciclo (Scrap View).
    \item L'Adapter, nel suo metodo \texttt{getView(position, convertView, parent)}, riceve come parametro \texttt{convertView} la view appena uscita dallo schermo.
    \item Se \texttt{convertView != null}, significa che possiamo \textbf{riutilizzare le View esistenti}.
    Non facciamo un nuovo "inflate" (che consumerebbe CPU e RAM), ma ci limitiamo a cambiare i testi/immagini al suo interno. Questo riduce il consumo di memoria e migliora vertiginosamente le performance.
\end{enumerate}

\begin{center}
        % Scala il grafico per riempire la larghezza del testo (\textwidth)
    \resizebox{\textwidth}{!}{
        \begin{tikzpicture}[
            node distance=1.5cm,
            box/.style={rectangle, draw, rounded corners, fill=white, text width=4cm, align=center, drop shadow},
            arrow/.style={-Stealth, thick, blue}
        ]
            \node[box, fill=orange!10] (getview) {\textbf{Adapter.getView()}\\Chiesto elemento \#10};
            \node[box, right=of getview, fill=gray!10] (check) {\texttt{convertView == null}?};
            \node[box, top color=red!10, bottom color=red!20, above right=of check] (inflate) {\textbf{SI:} LayoutInflate\\(Costoso!)};
            \node[box, top color=green!10, bottom color=green!20, below right=of check] (reuse) {\textbf{NO:} Riutilizzo View\\(Riciclo, Molto veloce!)};
            \node[box, fill=blue!10, right=4cm of check] (bind) {Associazione Dati\\(Set text/image)};

            \draw[arrow] (getview) -- (check);
            \draw[arrow] (check) |- (inflate);
            \draw[arrow] (check) |- (reuse);
            \draw[arrow] (inflate) -| (bind);
            \draw[arrow] (reuse) -| (bind);
        \end{tikzpicture}
    }
\end{center}

\subsection{SimpleAdapter vs ArrayAdapter (Domanda Quiz)}
Mentre l'\texttt{ArrayAdapter} viene utilizzato principalmente quando i dati sono contenuti in una
\textbf{lista semplice di oggetti} o array, il \textbf{\texttt{SimpleAdapter}} si usa per dati più complessi.
\begin{itemize}
    \item \texttt{SimpleAdapter} è più flessibile perché \textbf{può associare più campi a diverse View} (es. un titolo, un sottotitolo e un'immagine nella stessa riga).
    \item Viene usato quando i dati sono \textbf{strutturati come Map (chiave-valore)}, tipicamente in una \texttt{List<Map<String, Any>>}.
    \item L'array di stringhe fornito nel suo costruttore rappresenta \textbf{le chiavi della Map}, mentre l'array di interi rappresenta \textbf{gli ID delle View nel layout} in cui inserire i dati.
    \item Se si tenta di caricare un'immagine da un URL web tramite SimpleAdapter, l'immagine \textbf{non viene caricata automaticamente}.
    È necessario implementare un \texttt{ViewBinder} personalizzato e usare librerie come Glide.
\end{itemize}

\subsection{L'evoluzione: RecyclerView}
Essendo la \texttt{ListView} limitata, Google ha introdotto la \textbf{\texttt{RecyclerView}}.
\begin{itemize}
    \item \textbf{Gestione del Layout:} Delega il posizionamento a un \texttt{LayoutManager} (es. \texttt{LinearLayoutManager} per liste, \texttt{GridLayoutManager} per griglie).
    \item \textbf{Performance:} Impone (forza) per architettura l'uso del pattern \textbf{ViewHolder}, obbligando lo sviluppatore a mantenere i riferimenti in memoria degli elementi della riga, rendendo il riciclo estremamente efficiente.
\end{itemize}

\section{Jetpack Compose: Il paradigma moderno}
Oggi, Google spinge \textbf{Jetpack Compose}, un toolkit UI \textbf{dichiarativo} interamente scritto in Kotlin.
Non ci sono più file XML, non c'è più `findViewById`, né si usano gli Adapter!

In Compose:
\begin{itemize}
    \item \textbf{Tutto è una funzione:} Le UI si creano scrivendo funzioni annotate con \texttt{@Composable}.
    I nomi di queste funzioni si scrivono in PascalCase (es. \texttt{fun MyButton()}).
    \item \textbf{UI come funzione dello Stato:} \texttt{UI = f(state)}.
    Se lo stato (i dati) cambia, Compose effettua automaticamente la \textbf{Ricomposizione} (\textit{Recomposition}), ridisegnando solo le parti di UI colpite dal cambiamento.
    \item \textbf{Gestione di liste lunghe:} Al posto di \texttt{RecyclerView}, si usano \texttt{LazyColumn} (lista verticale) e \texttt{LazyRow} (orizzontale).
    Queste funzioni creano componenti solo quando diventano visibili sullo schermo, garantendo le stesse performance della RecyclerView, ma in 3 righe di codice.
\end{itemize}

\begin{tcolorbox}[colback=yellow!10!white,colframe=yellow!70!black,title=Confronto XML vs Compose (Utile per i quiz)]
    \begin{itemize}
        \item \texttt{LinearLayout (vertical)} $\rightarrow$ \texttt{Column}
        \item \texttt{LinearLayout (horizontal)} $\rightarrow$ \texttt{Row}
        \item \texttt{FrameLayout} $\rightarrow$ \texttt{Box}
        \item \texttt{RecyclerView} $\rightarrow$ \texttt{LazyColumn} / \texttt{LazyRow}
        \item \texttt{ListView.setOnItemClickListener} $\rightarrow$ \texttt{Modifier.clickable \{ \}}
    \end{itemize}
\end{tcolorbox}

\begin{lstlisting}[caption={Creare una lista in Jetpack Compose}]
@Composable
fun ListaUtenti(utenti: List<String>) {
    // LazyColumn equivale a una RecyclerView efficiente
    LazyColumn {
        // Il costrutto 'items' sostituisce l'Adapter!
        items(utenti) { utente ->
            Text(
                text = utente,
                modifier = Modifier
                            .fillMaxWidth()
                            .padding(16.dp)
            )
        }
    }
}
\end{lstlisting}

\section{Caricamento di Immagini da Rete}
Sia con XML che con Compose, scaricare immagini da Internet sul \textit{Main Thread} farebbe bloccare o crashare l'app.
Nelle app in Compose moderne, al posto di gestire thread separati, si utilizza la libreria \textbf{Coil} (Coroutine Image Loader), usando direttamente il widget \texttt{AsyncImage}:
\begin{lstlisting}[caption={AsyncImage con Coil}]
AsyncImage(
    model = "https://picsum.photos/200/300",
    contentDescription = "Immagine di prova",
    modifier = Modifier.size(64.dp)
)
\end{lstlisting}
\texttt{AsyncImage} esegue il download in modo asincrono senza bloccare l'interfaccia, salva l'immagine in cache e la mostra appena pronta.
Ricorda che è sempre necessario il permesso \texttt{INTERNET} nell' \texttt{AndroidManifest.xml}.